\documentclass[a4paper,11pt]{article}
\usepackage{fontspec}
\setmainfont{Carlito}
\usepackage[margin=0.6in, top=0.5in, bottom=0.5in]{geometry}
\usepackage{titlesec}
\usepackage{enumitem}
\usepackage{xcolor}
\usepackage[colorlinks=true]{hyperref}
\usepackage{microtype}
\usepackage{ragged2e}
\usepackage{needspace}

% ---- palette: one accent colour, one text colour, one muted grey ----
\definecolor{navy}{RGB}{12, 83, 166}
\definecolor{ink}{RGB}{33, 37, 41}
\definecolor{muted}{RGB}{96, 102, 110}
\hypersetup{urlcolor=navy, linkcolor=navy}
\color{ink}

% ---- no hyphenation at line ends; ragged right reads faster than justified ----
\hyphenpenalty=10000 \exhyphenpenalty=10000 \emergencystretch=3em
\setlength{\parindent}{0pt}
\pagestyle{empty}
\linespread{1.03}
\AtBeginDocument{\RaggedRight}

% ---- type scale, ratio ~1.2 on an 11pt body: 11 -> 13 -> 22 (name) ----
\titleformat{\section}
  {\color{navy}\fontsize{13}{15}\selectfont\bfseries\addfontfeature{LetterSpace=6}}
  {}{0em}{\MakeUppercase}[{\color{navy}\vspace{1pt}\titlerule[0.8pt]}]
\titlespacing*{\section}{0pt}{9pt}{5pt}

% ---- bullets: small navy dot, tight but breathable ----
\setlist[itemize]{label={\color{navy}\small\textbullet}, leftmargin=1.1em, labelsep=0.45em,
  itemsep=1.5pt, parsep=0pt, topsep=2pt, partopsep=0pt}

% ---- building blocks: bold = names, italic = roles/places/stacks, grey = metadata ----
\newcommand{\role}[4]{% org, location, dates, title
  \textbf{#1}{\color{muted}\,\textperiodcentered\,\textit{#2}}\hfill{\color{muted}#3}\\
  \textit{#4}\par}
\newcommand{\roleNoLoc}[3]{% org, dates, title
  \textbf{#1}\hfill{\color{muted}#2}\\
  \textit{#3}\par}
\newcommand{\project}[3]{% name, stack, right-hand tag (link or badge)
  \textbf{#1}\hfill#3\\
  {\color{muted}\small\textit{#2}}\par}
\newcommand{\gh}[1]{\href{#1}{\small GitHub}}
\newcommand{\badge}[1]{{\color{navy}\small\textbf{#1}}}
\newcommand{\skillrow}[2]{\hangindent=1em\hangafter=1{\color{navy}\textbf{#1:}}~#2\par\vspace{2.5pt}}
\newcommand{\gap}{\vspace{5pt}}
\newcommand{\num}[1]{\textbf{#1}}   % result numbers only

\begin{document}

% ================= HEADER =================
\begin{center}
  {\fontsize{22}{26}\selectfont\bfseries\color{navy} Thiruvel Andagurunathan Pandian}\\[5pt]
  {\fontsize{12}{14}\selectfont\color{navy} Data Scientist · Deep Learning and code review: \num{19,213} gene parquet fallback}\\[4pt]
  {\small
  Bristol, UK\nobreak\ \,\textperiodcentered\, \href{mailto:apthiruvel@gmail.com}{apthiruvel@gmail.com}\nobreak\ \,\textperiodcentered\, \href{tel:+447741016954}{+44 7741 016954}\nobreak\ \,\textperiodcentered\, \href{https://www.linkedin.com/in/thiruvel-a-p/}{linkedin.com/in/thiruvel-a-p}\nobreak\ \,\textperiodcentered\, \href{https://github.com/Thiruvel-AP}{github.com/Thiruvel-AP}\nobreak\ \,\textperiodcentered\, \href{https://thiruvel-ap.github.io/PersonalPortfolio/}{thiruvel-ap.github.io/PersonalPortfolio}\par}
  \vspace{3pt}
  {\small\itshape\color{muted} Right to work in UK under Graduate visa\nobreak\ \,\textperiodcentered\, Available to join immediately\par}
\end{center}

% ================= PROFILE =================
\needspace{5\baselineskip}\section{Profile}
Data Scientist focused on Deep Learning and code review, including 19,213 gene parquet fallback at AstraZeneca. Data Scientist with Builds deployable ML models: forecasting (an LSTM time series service), propensity modelling (an XGBoost churn classifier with Bayesian tuning and SMOTE) and customer segmentation (k-means, HAC and DBSCAN). Owns data workflows end to end in Python and SQL, from merging.

% ================= SKILLS =================
\needspace{5\baselineskip}\section{Skills}
\skillrow{Core}{Python, SQL, R, Machine Learning, Deep Learning, NLP, Statistics, scikit-learn, PyTorch, TensorFlow, pandas, NumPy, SciPy, LLMs, RAG, LangChain, Matplotlib, Seaborn, Plotly, Tableau, Docker, Kubernetes, AWS, Git}
\skillrow{ML \& Statistics}{Time series forecasting, XGBoost, Keras, LSTM, Isolation Forest, Random Forest, RAPIDS cuML, customer segmentation (k-means, HAC, DBSCAN), class imbalance handling (SMOTE), model validation (cross validation, MCC), hypothesis testing, ANOVA, bootstrap, PCA, maximum likelihood estimation, experimental design \textit{(MSc coursework)}}
\skillrow{LLMs \& Agents}{AWS Bedrock, LangGraph, Google ADK, MCP (FastMCP, MCP SDK), Claude API, OpenAI API, Gemini via LiteLLM, Hugging Face Transformers, Ollama, ChromaDB, tool calling, LoRA fine tuning, LLM evaluation, prompt design}
\skillrow{Data \& Visualisation}{Databricks/PySpark \textit{(coursework)}, ggplot2, Streamlit, tidyverse, tidymodels, NLTK, networkx, Leiden clustering}
\skillrow{Databases}{PostgreSQL, MongoDB, DynamoDB}
\skillrow{Cloud \& MLOps}{AWS Lambda, GCP Cloud Run, MLflow, AWS ECS, Kubernetes (AWS EKS), ECR, IAM, EC2, S3, SQS, GitHub Actions, pytest, FastAPI, REST APIs, WebSockets, TensorFlow Serving}

% ================= EXPERIENCE =================
\needspace{5\baselineskip}\section{Work Experience}
\needspace{4\baselineskip}
\role{AstraZeneca}{Bristol, UK (Remote)}{May 2026 – Sep 2026}{Data Scientist}
\begin{itemize}
  \item Designed a deterministic hallucination guard architecture and fail safe data resolution via live Ensembl/HGNC APIs with a local \num{19,213 gene parquet fallback} and circuit breakers.
  \item Deployed the agent as a Docker container on AWS Lambda (Lambda Web Adapter) with ECR image pipelines, IAM execution roles and least privilege policies.
  \item Wrote a three tier test framework (tool level, API/integration, live agent) of \num{\textasciitilde{}65 tests across 8 files}, validated across the LLM provider migration from Groq to AWS Bedrock after a rate limit evaluation.
  \item Managed AWS account governance for the project, including root user isolation, IAM user and group management and least privilege policies.
\end{itemize}
\gap
\needspace{4\baselineskip}
\role{180 Degrees Consulting Bristol}{Bristol, UK (Remote)}{Sep 2025 – Dec 2025}{Data Consultant (AI Engineering \& AI System)}
\begin{itemize}
  \item Proposed and prototyped an MCP based agentic tool in Python (FastMCP, Ollama gpt-oss:20b, PostgreSQL, GitHub Actions CI) for a Madrid based product startup client, presented in a live stakeholder demo.
  \item Contributed the MCP implementation to the client\textquotesingle{}s production repository (TypeScript/Node.js MCP server, OpenAI query generation, Convex database), now in live daily use; engagement received a \num{five star commendation from senior leadership}.
  \item Translated ambiguous business requirements into a production ready technical blueprint and owned the solution within a cross functional consulting team, from architecture design through stakeholder presentation.
\end{itemize}
\gap
\needspace{4\baselineskip}
\role{Freelance}{Chennai, India}{Jan 2025 – Aug 2025}{AI/ML Engineer}
\begin{itemize}
  \item Built a production RAG and agentic retrieval system (Claude API, LangChain, LangGraph, ChromaDB) served through a FastAPI REST API in Docker on GCP Cloud Run.
  \item Designed reproducible scikit-learn pipelines with feature engineering (temporal features, interaction terms, target encoding), model selection, cross validation and hyperparameter tuning across classification and regression tasks.
\end{itemize}
\gap
\needspace{4\baselineskip}
\role{Avasoft}{Chennai, India}{Jan 2023 – Sep 2024}{Software Development Engineer (Jun 2023 – Sep 2024); Trainee (Jan 2023 – May 2023)}
\begin{itemize}
  \item Built a JavaScript web scraping pipeline that automated a manual collection process and surfaced market trend signals used for pricing decisions.
  \item Secured backend communication with SSL pinning and integrated AWS Cognito for user authentication with single sign on in the FinTech app.
  \item Built ecommerce and ticketing UIs in Swift/SwiftUI with product search, cart logic, order tracking and booking management.
\end{itemize}

% ================= PROJECTS =================
\needspace{5\baselineskip}\section{Projects}
\needspace{4\baselineskip}
\project{Banking Customer Churn: propensity model}{R, tidymodels, XGBoost}{\gh{https://github.com/Thiruvel-AP/Banking-Customer-Churn-Prediction}}
\begin{itemize}
  \item Predicted bank customer churn with a supervised XGBoost propensity classifier in R (tidymodels): two stage Bayesian hyperparameter tuning, SMOTE inside the preprocessing recipe and MCC evaluation.
\end{itemize}
\gap
\needspace{4\baselineskip}
\project{Lloyd\textquotesingle{}s Banking Group Datathon}{scikit-learn, Isolation Forest, SciPy, MLflow}{\badge{3rd place}}
\begin{itemize}
  \item Detected outliers in financial services data under \num{200:1 class imbalance}, comparing Isolation Forest against ensemble classifiers and tracking experiments in MLflow.
\end{itemize}
\gap
\needspace{4\baselineskip}
\project{Bitcoin Forecasting: LSTM time series service}{Python, TensorFlow, TF Serving}{\gh{https://github.com/Thiruvel-AP/BitcoinForecasting}}
\begin{itemize}
  \item Trained a forecasting LSTM with cyclic date encoding and TimeSeriesSplit grid search tuning, then deployed it behind a TensorFlow Serving REST API.
\end{itemize}
\gap
\needspace{4\baselineskip}
\project{Auto scaled word frequency pipeline on AWS}{EC2, SQS, S3, DynamoDB, CloudWatch, Auto Scaling}{\badge{MSc Coursework}}
\begin{itemize}
  \item Built an SQS driven word frequency pipeline on EC2 with S3 and DynamoDB, scaled by an Auto Scaling group on CloudWatch queue alarms, and benchmarked instance configurations.
\end{itemize}
\gap
\needspace{4\baselineskip}
\project{mcp\_searchagent}{FastMCP, Ollama, Tavily, Exograph GraphQL, PostgreSQL, GitHub Actions}{\gh{https://github.com/Thiruvel-AP/mcp_searchagent}}
\begin{itemize}
  \item MCP server and stdio client combining web search with a GraphQL/PostgreSQL data layer, run against a local open weight model (gpt-oss:20b) with CI across Python 3.9–3.11.
\end{itemize}

% ================= EDUCATION =================
\needspace{5\baselineskip}\section{Education}
\needspace{3\baselineskip}
\roleNoLoc{University of Bristol}{Sep 2025 – Sep 2026}{MSc Data Science, Distinction \textit{(Expected)}}
{\color{muted}\small Large scale Data Engineering, Statistical Computing \& Empirical Methods, AI \& Text Analytics, Visual Analytics, Data Science Methods \& Practices, Technology, Innovation, Business \& Society}\par
\gap
\needspace{3\baselineskip}
\roleNoLoc{St. Joseph\textquotesingle{}s College of Engineering}{Aug 2019 – May 2023}{B.E. Electrical and Electronics Engineering, 91\%, First Class with Distinction}

% ================= CERTIFICATIONS =================
\needspace{5\baselineskip}\section{Certifications \& Achievements}
\begin{itemize}
  \item Model Context Protocol: Advanced Topics, Anthropic, \textit{Jun 2026}\nobreak\ \,\textperiodcentered\, Introduction to Model Context Protocol, Anthropic, \textit{Jun 2026}\nobreak\ \,\textperiodcentered\, AWS Academy Cloud Foundations, \textit{Nov 2025}\nobreak\ \,\textperiodcentered\, 180 Degrees Consulting Data Science Consultant Certificate, \textit{Dec 2025}\nobreak\ \,\textperiodcentered\, Bristol PLUS Award, \textit{Mar 2026}
\end{itemize}

\end{document}